%% Machine Learning and Agentic AI for Macroeconomic Research
%% Lecture 3: Machine Learning for Quantitative Macroeconomists
%% Author: Zhigang Feng
%% Source: 5-machine_learning/Lec_2026_5_Machine_learning.tex

\documentclass[10pt,english,aspectratio=169]{beamer}

\usepackage{ae,aecompl}
\usepackage[T1]{fontenc}
\usepackage{textcomp}
\usepackage[utf8]{inputenc}
\usepackage{babel}

\setcounter{secnumdepth}{3}
\setcounter{tocdepth}{3}

\usepackage{amsmath, amssymb, amsthm, graphicx}
\usepackage{booktabs, multirow}
\usepackage{tcolorbox}
\tcbuselibrary{skins, breakable, theorems}
\usepackage{listings}
\usepackage{xcolor}

\lstset{
  basicstyle=\ttfamily\small,
  keywordstyle=\color{blue!80!black}\bfseries,
  commentstyle=\color{green!50!black}\itshape,
  stringstyle=\color{orange!80!black},
  backgroundcolor=\color{gray!8},
  frame=single, rulecolor=\color{gray!40},
  breaklines=true, showstringspaces=false,
  numbers=left,
  numbersep=5pt,
  numberstyle=\tiny\color{gray},
}

\lstdefinestyle{pythonstyle}{
    language=Python,
    basicstyle=\ttfamily\footnotesize,
    keywordstyle=\color{blue!80!black}\bfseries,
    commentstyle=\color{green!50!black}\itshape,
    stringstyle=\color{orange!80!black},
    frame=single,
    breaklines=true,
    showstringspaces=false,
    numbers=left,
    numbersep=5pt,
    numberstyle=\tiny\color{gray},
    backgroundcolor=\color{gray!8},
    rulecolor=\color{gray!40}
}

\ifx\hypersetup\undefined
  \AtBeginDocument{%
    \hypersetup{bookmarks=false,breaklinks=true,pdfborder={0 0 0},
      colorlinks=true,linkcolor=DarkText,urlcolor=RedTitle}%
  }
\else
  \hypersetup{bookmarks=false,breaklinks=true,pdfborder={0 0 0},
    colorlinks=true,linkcolor=DarkText,urlcolor=RedTitle}%
\fi

\makeatletter
\newcommand\makebeamertitle{\frame{\maketitle}}%
\AtBeginDocument{%
  \let\origtableofcontents=\tableofcontents
  \def\tableofcontents{\@ifnextchar[{\origtableofcontents}{\gobbletableofcontents}}
  \def\gobbletableofcontents#1{\origtableofcontents}%
}

\usetheme{metropolis}
\usefonttheme{professionalfonts}

\definecolor{LightGray}{RGB}{230,230,230}
\definecolor{RedTitle}{RGB}{0,0,200}
\definecolor{VeryLightGray}{RGB}{245,245,245}
\definecolor{DarkText}{RGB}{0,0,0}

\setbeamercolor{background canvas}{bg=white}
\setbeamercolor{title}{fg=RedTitle}
\setbeamercolor{author}{fg=DarkText}
\setbeamercolor{date}{fg=DarkText}
\setbeamercolor{frametitle}{bg=VeryLightGray,fg=DarkText}
\setbeamercolor{normal text}{fg=DarkText}

\setbeamersize{text margin left=5mm, text margin right=5mm}
\addtobeamertemplate{frametitle}{}{\vspace{-0.5em}}
\newcommand{\reducevspace}{\vspace{-0.5cm}}

\setbeamertemplate{navigation symbols}{}
\setbeamertemplate{footline}{%
  \hfill%
  \usebeamercolor[fg]{page number in head/foot}%
  \usebeamerfont{page number in head/foot}%
  \insertframenumber\,/\,\inserttotalframenumber\kern1em\vskip2pt%
}

\makeatother

\graphicspath{{../../5-machine_learning/}}

\title[ML \& Agentic AI for Macro Research]{Machine Learning and Agentic AI for Macroeconomic Research\\
3: Machine Learning for Quantitative Macroeconomists\\[0.5em]}
\author{Zhigang Feng}
\date{Spring 2026}

\begin{document}

\makebeamertitle

%=============================================================================
% AGENDA
%=============================================================================

\begin{frame}{Lecture Agenda}
\begin{enumerate}
  \item \textbf{Why ML for Economists?} --- unstructured data, forecasting, solving macro models
  \medskip
  \item \textbf{Pattern Recognition} --- handwritten digits as a motivating illustration
  \medskip
  \item \textbf{Neural Networks} --- architecture, activation functions, universal approximation
  \medskip
  \item \textbf{Forward Propagation} --- computing outputs layer by layer
  \medskip
  \item \textbf{Backpropagation \& Gradient Descent} --- learning from errors
  \medskip
  \item \textbf{Training Process} --- loss functions, overfitting, regularization
  \medskip
  \item \textbf{Specialized Architectures} --- CNNs, convexity- \& monotonicity-preserving networks
  \medskip
  \item \textbf{Lab Preview} --- curve fitting with \texttt{Lab5A} and \texttt{Lab5B}
\end{enumerate}
\end{frame}

%=============================================================================
% SECTION 1: WHY ML FOR ECONOMISTS
%=============================================================================
\section{Why ML for Economists?}

\begin{frame}{Harnessing the Power of Unstructured Data}
\begin{itemize}
\item Economics has traditionally relied on structured, tabular data. However, a wealth of economic information is locked in unstructured formats: text, images, satellite imagery.
\item ML algorithms, particularly deep learning, provide tools to systematically process and quantify this unstructured data into economic indicators.
\item \textbf{Applications:}
\begin{itemize}
\item Measuring economic sentiment and uncertainty from financial news and social media
\item Analyzing satellite images to track economic activity (nighttime lights $\approx$ GDP proxy, ships in ports $\approx$ trade flows)
\item Using image recognition on product pictures to understand non-standardized product features
\end{itemize}
\end{itemize}
\end{frame}

\begin{frame}{Extracting Insights from Textual Data}
\begin{itemize}
\item Natural Language Processing (NLP) offers powerful techniques to analyze large volumes of text --- going beyond simple word counts.
\item \textbf{Applications:}
\begin{itemize}
\item Gauging dovish or hawkish stance of central bankers from public communications
\item Identifying key risk factors from corporate annual reports
\item Quantifying the evolution of economic policy uncertainty from news articles
\end{itemize}

\medskip
\item \textbf{For macro:} FOMC minutes, Beige Books, and congressional records contain rich information that traditional econometrics cannot easily exploit.
\end{itemize}
\end{frame}

\begin{frame}{Enhancing Causal Inference and Forecasting}
\begin{itemize}
\item \textbf{Causal Inference:}
\begin{itemize}
\item ``Double Machine Learning'' allows flexible inclusion of high-dimensional controls, reducing omitted variable bias
\item ML identifies \emph{heterogeneous treatment effects}: how does a policy impact vary across subgroups?
\end{itemize}

\medskip
\item \textbf{Forecasting \& Nowcasting:}
\begin{itemize}
\item ML captures complex, non-linear relationships that time-series models miss
\item Regularization (e.g., LASSO) enables ``kitchen sink'' regression without overfitting
\item Nowcasting: real-time economic indicators from credit card transactions, web searches, satellite imagery
\end{itemize}
\end{itemize}
\end{frame}

\begin{frame}{Solving Large-Scale Macro and Finance Models}
\begin{itemize}
\item Cutting-edge models (heterogeneous agents, portfolio choice) feature:
\begin{itemize}
\item High dimensionality and complex non-linear dynamics
\item Traditional grid-based methods hit the \textbf{curse of dimensionality}
\end{itemize}

\medskip
\item ML techniques provide powerful new numerical methods:
\begin{itemize}
\item \textbf{Deep learning:} parameterize policy/value functions with neural networks
\item \textbf{Reinforcement learning:} solve sequential decision problems via simulation
\end{itemize}

\medskip
\item \textbf{Examples:}
\begin{itemize}
\item Optimal fiscal/monetary policy in heterogeneous agent models
\item Life-cycle portfolio choice with realistic income uncertainty
\item Krusell--Smith (1998) with neural network approximation of the law of motion
\end{itemize}
\end{itemize}
\end{frame}

\begin{frame}{ML and the Curse of Dimensionality in Macro}
\begin{itemize}
\item \textbf{CoD \#1: Dimension of the state space}
\begin{itemize}
\item $k\rightarrow\{k,a,h,\ldots\}$; the distribution $\Gamma$ is infinite-dimensional
\item \textit{Traditional:} sparse grids, local approximation --- scale poorly
\item \textit{ML:} neural networks are \textbf{grid-free}; recover equilibrium functions via Monte Carlo sampling
\end{itemize}

\medskip
\item \textbf{CoD \#2: Unfolding of uncertainty}
\begin{itemize}
\item $\mathbb{E}$ over future contingencies; the space for $\Gamma$ is unknown
\item \textit{Traditional:} discretize shocks, impose Markov structure
\item \textit{ML:} focus on the ergodic distribution; direct simulation Monte Carlo
\end{itemize}

\medskip
\item \textbf{Key insight:} ML shifts the computational bottleneck from \emph{grid construction} to \emph{function representation} --- a much more scalable paradigm.
\end{itemize}
\end{frame}

%=============================================================================
% SECTION 2: PATTERN RECOGNITION
%=============================================================================
\section{Pattern Recognition as Motivating Illustration}

\begin{frame}{Example: Handwritten Digit Recognition}
\begin{center}
\includegraphics[width=12cm,height=7cm]{3.png}
\end{center}
\end{frame}

\begin{frame}{Digit Recognition as a Mathematical Problem}
\begin{itemize}
  \item \textbf{Dataset:}
        $\bigl\{(\mathbf{x}_i, y_i)\bigr\}_{i=1}^N$ with images $\mathbf{x}_i \in [0,1]^{784}$
        and labels $y_i \in \{0,\dots,9\}$

  \item \textbf{Input--Output Mapping:}
        \[
          f_{\theta} \colon \mathbf{x} \longmapsto \hat{\mathbf{p}} \in \Delta^{9},
          \qquad
          \hat{y} = \arg\max_{k} \hat{p}_k
        \]
        where $\Delta^{9}$ is the 10-dimensional probability simplex

  \item \textbf{Learning Objective (Empirical Risk + Regularization):}
        \[
          \min_{\theta} \;
          \frac{1}{N}\sum_{i=1}^{N}
          \underbrace{\mathcal{L}\bigl(f_{\theta}(\mathbf{x}_i), y_i\bigr)}_{\text{cross-entropy}}
          + \lambda \lVert\theta\rVert_{p}
        \]

  \item \textbf{Search Strategy:} Stochastic gradient descent \& variants (Adam, RMSProp)
\end{itemize}
\end{frame}

\begin{frame}{Constraints and Structure in Handwritten Digits}
\begin{itemize}
\item \textbf{Anatomical constraints:} Muscles, tendons, and joints dictate possible strokes --- leading to systematic patterns in how digits are formed.

\medskip
\item \textbf{Geometric properties:} Each digit has defining features --- `0' has a closed loop, `1' a vertical stroke, `8' two stacked loops.

\medskip
\item \textbf{Spatial relationships:} Relative positions, sizes, and orientations of components maintain digit identity even with individual variation.

\medskip
\item \textbf{Lesson for economists:} Just as economic theory constrains the space of plausible models, structural constraints on handwriting reduce the effective dimensionality of the recognition problem.

\medskip
\item ML algorithms can \textbf{learn} these constraints from data --- data augmentation, regularization, and architecture choices all exploit this structure.
\end{itemize}
\end{frame}

\begin{frame}{From 784 Pixels to a Low-Dimensional Manifold}
\begin{itemize}
  \item \textbf{Manifold Hypothesis:}
        Natural digit images occupy a smooth, $d$-dimensional manifold
        $\mathcal{M}\subset\mathbb{R}^{784}$ with $d\ll784$.

  \item \textbf{Two-Stage Model:}
        \[
          \mathbf{x}
          \xrightarrow{\;\pi\,:\;\mathbb{R}^{784}\!\!\to\mathbb{R}^{d}\;}
          \mathbf{z}
          \xrightarrow{\;g_{\phi}\,:\;\mathbb{R}^{d}\!\!\to\Delta^{9}\;}
          \hat{\mathbf{p}}
          ,\qquad
          \hat y=\arg\max_k\hat p_k
        \]
        \begin{itemize}
          \item $\pi$: \textit{encoder} --- flattens the manifold to a latent code $\mathbf z\in\mathbb{R}^{d}$
          \item $g_{\phi}$: \textit{classifier head} --- 1--2 FC layers + softmax $\to$ class probabilities
        \end{itemize}

  \item \textbf{Choices for $\pi$:}
        \begin{itemize}
          \item \textbf{Linear:} PCA, SVD, random projections
          \item \textbf{Non-linear:} Auto-encoder, t-SNE, UMAP, first layers of a CNN
        \end{itemize}

  \item \textbf{Macro analogy:} Finding a low-dimensional state in a high-dimensional model is the same problem --- sufficient statistics for prediction.
\end{itemize}
\end{frame}

\begin{frame}{Defining a Deep Neural Network for Digit Recognition}
\begin{itemize}
\item \textbf{Network Architecture:}
\begin{itemize}
\item Input Layer: 784 neurons corresponding to pixel intensities
\item Hidden Layers: Multiple layers applying nonlinear transformations to learn complex patterns
\item Output Layer: 10 neurons, each representing one digit class (0--9)
\end{itemize}

\medskip
\item \textbf{Activation Functions:}
\begin{itemize}
\item Nonlinear functions (ReLU, sigmoid) between layers
\item Customizable depth and width based on task complexity
\end{itemize}

\medskip
\item \textbf{Economic parallel:} The network structure mirrors the recursive structure in macro models --- state variables are transformed layer by layer into decisions, just as $k_t \to c_t, l_t$ through policy functions.
\end{itemize}
\end{frame}

%=============================================================================
% SECTION 3: NEURAL NETWORKS
%=============================================================================
\section{Neural Networks: Architecture and Activation Functions}

\begin{frame}{Neural Networks: The Basic Building Block}
\begin{itemize}
\item Artificial neural networks are:
\begin{itemize}
\item Inspired by biological neural networks
\item Composed of interconnected neurons organized in layers
\item Trained by adjusting weights based on data
\end{itemize}

\medskip
\item \textbf{Key components:}
\begin{itemize}
\item \textbf{Neuron:} Receives inputs, applies an activation function, produces output
\item \textbf{Activation function:} Nonlinear map applied by each neuron (sigmoid, tanh, ReLU)
\item \textbf{Loss function:} Measures how well predictions match true values --- the objective to minimize
\end{itemize}
\end{itemize}

\begin{center}
\includegraphics[scale=0.5]{neuron.png}
\end{center}
\end{frame}

\begin{frame}{Single Hidden Layer Network}
\begin{itemize}
\item Input $x\in\mathbb{R}^{M}$, output $y\in\mathbb{R}^{N}$.

\item The $j$-th hidden unit:
\begin{equation}
h_{j}^{(1)} = f_{j}^{(1)}(x) = \sigma\!\left(b_{j}^{(1)}+\sum_{m=1}^{M}w_{j,m}^{(1)}x_{m}\right),\quad j=1,\dots,U_{1} \label{eq:hidden1}
\end{equation}

\item The $i$-th output:
\begin{equation}
y_{i} = \sum_{j=1}^{U_{1}}\tilde{w}_{i,j}\,h_{j}^{(1)}+\tilde{b}_{i},\quad i=1,\dots,N \label{eq:output1}
\end{equation}

\item $\sigma$ is a nonlinear activation function.

\medskip
\item Common choices: $\sigma(x)=\tanh(x)$, $\sigma(x)=\max(0,x)$ (ReLU), $\sigma(x)=\frac{1}{1+e^{-x}}$ (sigmoid).
\end{itemize}
\end{frame}

\begin{frame}{Two Hidden Layer Network}
\begin{itemize}
\item Take the linear combination of $h^{(1)}$ and pass through another activation:
\begin{equation}
h_{j}^{(2)} = \sigma\!\left(b_{j}^{(2)}+\sum_{k=1}^{U_{1}}w_{j,k}^{(2)}\,f_{k}^{(1)}(x)\right),\quad j=1,\dots,U_{2} \label{eq:hidden2}
\end{equation}
\begin{equation}
y_{i} = \sum_{j=1}^{U_{2}}\tilde{w}_{i,j}\,h_{j}^{(2)}+\tilde{b}_{i},\quad i=1,\dots,N \label{eq:output2}
\end{equation}

\medskip
\item \textbf{Pattern:} Each layer performs a linear transformation followed by a nonlinear activation. Depth adds representational power.

\medskip
\item \textbf{Macro analogy:} Successive layers are like successive time steps in a dynamic program --- each transforms the ``state'' into a richer representation.
\end{itemize}
\end{frame}

\begin{frame}{Hidden Layers: Visual}
\begin{center}
\includegraphics[width=12cm]{dnn}
\end{center}
\end{frame}

\begin{frame}{Activation Functions}
\begin{center}
\includegraphics[width=12cm]{activation}
\end{center}

\begin{itemize}
\item \textbf{Sigmoid:} $\sigma(x) = \frac{1}{1+e^{-x}}$ --- outputs in $(0,1)$, used for probabilities
\item \textbf{Tanh:} $\tanh(x)$ --- outputs in $(-1,1)$, zero-centered
\item \textbf{ReLU:} $\max(0,x)$ --- simple, sparse activation, dominant in modern networks
\end{itemize}
\end{frame}

\begin{frame}{Function Approximation with ReLU}
\begin{itemize}
    \item \textbf{Goal:} Approximate arbitrary nonlinear functions using simple building blocks.
    \item \textbf{Key Idea:} Linear combination of shifted ReLU functions:
    \[
    \hat{y}(x) = a_0 + \sum_{i=1}^{N} a_i\,\text{ReLU}(x-b_i)
    \]
    \item \textbf{Shifts $b_i$:} Define the ``kink'' points where marginal slope changes.
    \item \textbf{Coefficients $a_i$:} Scale each ReLU's contribution to match local slope.

    \medskip
    \item \textbf{Example:} Approximating $\sin(x)$ on $[-\pi, \pi]$:
    \begin{itemize}
        \item Choose shifts $b_i$ evenly over the domain
        \item Form design matrix with columns $\text{ReLU}(x-b_i)$ and a bias
        \item Solve a least-squares problem for coefficients $a_i$
    \end{itemize}

    \medskip
    \item \textbf{Economist's take:} ReLU kinks are piecewise-linear splines --- the network \emph{learns} the knot locations, unlike a fixed Chebyshev grid.
\end{itemize}
\end{frame}

\begin{frame}{ReLU Approximation: Increasing Complexity}
\begin{center}
\includegraphics[width=12cm,height=7.5cm]{relu1}
\end{center}
\end{frame}

\begin{frame}{ReLU Approximation: More Kink Points}
\begin{center}
\includegraphics[width=12cm,height=7.5cm]{relu2}
\end{center}
\end{frame}

\begin{frame}{ReLU Approximation: High Fidelity}
\begin{center}
\includegraphics[width=12cm,height=7.5cm]{relu3}
\end{center}
\end{frame}

\begin{frame}{Why Neural Networks Work: Theoretical Guarantees}

\begin{tcolorbox}[colback=blue!3, colframe=blue!40, title={Universal Approximation Theorem (Cybenko, 1989)}]
If $\sigma$ is sigmoidal (i.e., $\lim_{z\to-\infty}\sigma(z)=0$, $\lim_{z\to+\infty}\sigma(z)=1$), then any continuous function on $[0,1]^d$ can be approximated uniformly by a two-layer neural network.
\end{tcolorbox}

\medskip

\begin{tcolorbox}[colback=green!3, colframe=green!40, title={Breaking the Curse of Dimensionality (Barron, 1993)}]
A one-layer neural network achieves integrated squared error of order $O(1/M)$, where $M$ is the number of nodes. In contrast, series approximations (polynomials, splines) achieve $O(1/M^{2/N})$ where $N$ is the input dimension --- exponentially worse.
\end{tcolorbox}

\medskip
\begin{itemize}
\item \textbf{For economists:} Barron's result explains why neural networks scale to high-dimensional state spaces where Chebyshev polynomials break down.
\end{itemize}
\end{frame}

\begin{frame}{DNN = Linear Transformation + Nonlinear Activation}
\begin{itemize}
\item \textbf{Linear Transformation:}
\begin{itemize}
\item Each neuron performs a weighted sum of inputs: projects data onto a new axis
\item The set of transformations across a layer creates a \emph{new representation}
\item Enables: flexible feature representation, information compression, interaction capture
\end{itemize}

\medskip
\item \textbf{Nonlinear Activation:}
\begin{itemize}
\item Allows the network to capture nonlinearity and complex patterns
\item Creates curved decision boundaries in the feature space
\item \textbf{Essential} for the universal approximation property
\end{itemize}

\medskip
\item \textbf{Together:} Each layer applies an affine map $Wx+b$ then a pointwise nonlinearity $\sigma(\cdot)$. Stacking layers composes these maps:
\[
f(x) = \sigma_L \circ A_L \circ \cdots \circ \sigma_1 \circ A_1 (x)
\]
where $A_\ell(x) = W_\ell x + b_\ell$.
\end{itemize}
\end{frame}

%=============================================================================
% SECTION 4: FORWARD PROPAGATION
%=============================================================================
\section{Forward Propagation}

\begin{frame}{Feedforward: Computing Outputs Layer by Layer}
\begin{itemize}
\item \textbf{Forward propagation} (feedforward): compute the network output given an input by passing through layers sequentially.

\medskip
\item For a network with $L$ hidden layers:
\[
\mathbf{h}^{(0)} = \mathbf{x}, \qquad
\mathbf{h}^{(\ell)} = \sigma\!\left(W^{(\ell)}\mathbf{h}^{(\ell-1)} + \mathbf{b}^{(\ell)}\right), \quad \ell = 1,\ldots,L
\]
\[
\hat{\mathbf{y}} = W^{(L+1)}\mathbf{h}^{(L)} + \mathbf{b}^{(L+1)}
\]

\medskip
\item Each layer: (1) multiply by weight matrix, (2) add bias, (3) apply activation.

\medskip
\item \textbf{Analogy:} Forward propagation is like solving a recursive model forward in time --- given initial state $x$, compute successive transformations to arrive at the output.
\end{itemize}
\end{frame}

\begin{frame}{Forward Pass: Visual}
\begin{center}
\includegraphics[width=12cm]{nn_working1}
\end{center}
\end{frame}

\begin{frame}{Toy Example: Forward Pass with Numbers}
Consider a tiny network: input $x=1.0$, one hidden neuron (sigmoid), one output neuron (sigmoid).

\medskip
Parameters: $w_1 = 0.5,\; b_1 = 0,\; w_2 = 2.0,\; b_2 = 0$

\medskip
\textbf{Step 1: Hidden layer}
\[
z_1 = w_1 \cdot x + b_1 = 0.5 \quad\Rightarrow\quad h = \sigma(0.5) \approx 0.62
\]

\textbf{Step 2: Output layer}
\[
z_2 = w_2 \cdot h + b_2 = 2.0 \times 0.62 \approx 1.24 \quad\Rightarrow\quad \hat{y} = \sigma(1.24) \approx 0.78
\]

\textbf{Step 3: Loss} (target $t=0.0$)
\[
L = \tfrac{1}{2}(\hat{y} - t)^2 = \tfrac{1}{2}(0.78)^2 \approx 0.30
\]

\medskip
This is a complete forward pass: from $x$ to $\hat{y}$ to loss $L$.
\end{frame}

%=============================================================================
% SECTION 5: BACKPROPAGATION AND GRADIENT DESCENT
%=============================================================================
\section{Backpropagation and Gradient Descent}

\begin{frame}{Backpropagation: Learning from Errors}
\begin{itemize}
\item \textbf{Goal:} Compute $\frac{\partial L}{\partial \theta}$ for every parameter $\theta$ in the network.

\medskip
\item \textbf{Method:} Apply the \textbf{chain rule} backwards through the network layers.

\medskip
\item \textbf{Steps:}
\begin{enumerate}
\item Forward pass: compute all activations and the loss
\item Backward pass: propagate error signal from output to input
\item Compute gradient of loss w.r.t.\ each weight and bias
\item Update parameters: $\theta \leftarrow \theta - \eta \cdot \nabla_\theta L$
\end{enumerate}
\end{itemize}

\begin{center}
\includegraphics[width=10cm]{nn_working3}
\end{center}
\end{frame}

\begin{frame}{Toy Example: Backpropagation --- Output Layer}

Continuing our numerical example ($\hat{y} \approx 0.78$, $h \approx 0.62$, target $t=0$):

\medskip
\textbf{Output layer gradients:}
\[
\frac{\partial L}{\partial \hat{y}} = \hat{y} - t \approx 0.78
\]
\[
\frac{\partial \hat{y}}{\partial z_2} = \hat{y}(1-\hat{y})
\approx 0.78 \times 0.22 \approx 0.17
\]
\[
\frac{\partial L}{\partial z_2}
= \frac{\partial L}{\partial \hat{y}} \cdot \frac{\partial \hat{y}}{\partial z_2}
\approx 0.78 \times 0.17 \approx 0.13
\]
\[
\frac{\partial L}{\partial w_2}
= \frac{\partial L}{\partial z_2} \cdot \frac{\partial z_2}{\partial w_2}
= 0.13 \times h \approx 0.13 \times 0.62 \approx 0.08
\]

\medskip
\textbf{Interpretation:} The gradient $0.08$ tells us: to reduce the loss, decrease $w_2$ (the output overshot the target of $0$).
\end{frame}

\begin{frame}{Toy Example: Backpropagation --- Hidden Layer \& Update}

\textbf{Hidden layer gradients:}
\[
\frac{\partial L}{\partial h}
= \frac{\partial L}{\partial z_2} \cdot w_2
= 0.13 \times 2.0 \approx 0.26
\]
\[
\frac{\partial h}{\partial z_1} = h(1-h)
\approx 0.62 \times 0.38 \approx 0.24
\]
\[
\frac{\partial L}{\partial w_1}
= \frac{\partial L}{\partial h} \cdot \frac{\partial h}{\partial z_1} \cdot x
= 0.26 \times 0.24 \times 1.0 \approx 0.06
\]

\medskip
\textbf{Parameter update} (learning rate $\eta = 0.1$):
\[
w_2^{\text{new}} \approx 2.0 - 0.1 \times 0.08 = 1.992, \qquad
w_1^{\text{new}} \approx 0.5 - 0.1 \times 0.06 = 0.494
\]

\medskip
\textbf{Key point:} This is exactly what PyTorch's \texttt{.backward()} and the optimizer do --- just at massive scale with millions of parameters.
\end{frame}

\begin{frame}{Gradient Descent: The Optimization Engine}

\begin{itemize}
\item \textbf{Core update rule:}
$\theta_{n+1} \leftarrow \theta_n - \lambda_n \, g(b_n; \theta_n)$

\medskip
\item \textbf{Variants by batch size $n$:}
\begin{itemize}
\item $n=1$: \textbf{Stochastic Gradient Descent (SGD)} --- one random sample per step
\item $n=N$: \textbf{Full Gradient Descent (GD)} --- all data per step
\item $1 < n < N$: \textbf{Mini-batch GD} --- the practical sweet spot
\end{itemize}

\medskip
\item \textbf{Stochastic gradient approximation:}
\[
\frac{d\mathcal{L}(\Gamma_{\gamma})}{d\gamma} \approx \frac{1}{n}\sum_{i=1}^{n} \frac{d\mathcal{L}(\Gamma_{\gamma}; b_i)}{d\gamma}
\]

\medskip
\item \textbf{Algorithm:}
\begin{enumerate}
\item Initialize parameters randomly
\item Draw random mini-batch $b_n$; compute stochastic gradient $g(b_n; \theta_n)$
\item Choose learning rate $\lambda_n > 0$; update $\theta_{n+1} \leftarrow \theta_n - \lambda_n \, g(b_n; \theta_n)$
\item Repeat until convergence
\end{enumerate}
\end{itemize}
\end{frame}

\begin{frame}{Gradient Descent: Visual}
\begin{center}
\includegraphics[width=12cm]{nn_working6}
\end{center}
\end{frame}

\begin{frame}{Adaptive Learning Rates: RMSProp and Adam}
\begin{itemize}
\item \textbf{RMSProp:} Adjust learning rates individually by parameter:
\begin{align*}
g_{k} &\leftarrow \frac{1}{n}\sum_{i=1}^{n}\nabla\mathcal{L}(\Gamma_{\gamma}; b_i) \\
r_{k+1} &\leftarrow \rho\, r_k + (1-\rho)\, g_k \odot g_k \\
\Delta\theta_{k+1} &= -\frac{\lambda}{\sqrt{\delta + r_{k+1}}} \odot g_k
\end{align*}
Divides by running average of squared gradients --- normalizes step size.

\medskip
\item \textbf{Adam} (Adaptive Moments): Combines momentum + RMSProp.
\begin{itemize}
\item Tracks exponential moving average of gradient (1st moment) and squared gradient (2nd moment)
\item Parameters $\rho_1, \rho_2$ control decay rates
\item Default choice in most deep learning applications
\end{itemize}
\end{itemize}
\end{frame}

\begin{frame}{Batch vs.\ Stochastic vs.\ Mini-Batch}
\begin{center}
\includegraphics[width=12cm]{batch_1}
\end{center}
{\small \textbf{Batch GD:} Smooth convergence path, but each step is expensive (uses all data).}
\end{frame}

\begin{frame}{Stochastic Gradient Descent}
\begin{center}
\includegraphics[width=12cm]{batch_2}
\end{center}
{\small \textbf{SGD:} Noisy path with high variance, but each step is cheap and can escape local minima.}
\end{frame}

\begin{frame}{Mini-Batch Gradient Descent}
\begin{center}
\includegraphics[width=12cm]{minibatch}
\end{center}
{\small \textbf{Mini-batch:} Balances speed (parallelizable) and stability (averaged gradients). The practical standard.}
\end{frame}

\begin{frame}{Why Mini-Batch? Four Reasons}
\begin{enumerate}
\item \textbf{Computational efficiency:} More frequent parameter updates than full GD; faster convergence in wall-clock time.

\medskip
\item \textbf{Memory constraints:} Large datasets may not fit in GPU memory. Mini-batches process manageable subsets.

\medskip
\item \textbf{Implicit regularization:} The stochasticity in gradient estimates helps escape local minima and saddle points --- acts as a form of regularization.

\medskip
\item \textbf{Parallelization:} Each mini-batch can be processed across multiple GPUs simultaneously.
\end{enumerate}

\medskip
\textbf{Economist's note:} Mini-batch SGD is the computational analogue of sampling-based estimation in econometrics --- trade exactness for speed, and the noise is actually helpful.
\end{frame}

\begin{frame}{Convergence Theory (For the Technically Inclined)}

\begin{tcolorbox}[colback=blue!3, colframe=blue!40, title=Convergence with Diminishing Learning Rate]
Under standard assumptions (Lipschitz gradient, bounded variance), if the learning rate sequence satisfies
\[
\sum_{n=1}^{N}\lambda_n = \infty, \qquad \sum_{n=1}^{N}\lambda_n^2 < \infty
\]
then the expected average squared gradient converges:
\[
\mathbb{E}\!\left[\frac{1}{\Delta_N}\sum_{n=1}^{N}\lambda_n \|\nabla F(\theta_n)\|^2\right] \xrightarrow{N\to\infty} 0
\]
where $\Delta_N = \sum_{n=1}^N \lambda_n$.

\medskip
{\small --- Bottou, Curtis, Nocedal (2018). \textit{Optimization Methods for Large-Scale ML}.}
\end{tcolorbox}

\medskip
\begin{itemize}
\item The conditions $\sum \lambda_n = \infty$ (keep exploring) and $\sum \lambda_n^2 < \infty$ (shrink steps) mirror Robbins--Monro conditions familiar from stochastic approximation in econometrics.
\end{itemize}
\end{frame}

%=============================================================================
% SECTION 6: TRAINING PROCESS
%=============================================================================
\section{Training Process: Loss Functions, Overfitting, Regularization}

\begin{frame}{Loss Functions}
\begin{itemize}
\item \textbf{Mean Squared Error (regression):}
\[
\mathcal{L}_{\text{MSE}} = \frac{1}{N}\sum_{i=1}^N (\hat{y}_i - y_i)^2
\]
Equivalent to maximum likelihood under Gaussian errors.

\medskip
\item \textbf{Cross-Entropy (classification):}
\[
\mathcal{L}_{\text{CE}} = -\sum_{k=0}^{K-1} y_k \log(\hat{y}_k)
\]
where $y_k$ is the one-hot ground truth and $\hat{y}_k = \text{softmax}(z)_k = \frac{\exp(z_k)}{\sum_j \exp(z_j)}$.

\medskip
\item \textbf{Intuition:} Cross-entropy is small when confident \emph{and correct}, large when confident \emph{but wrong}.

\medskip
\item \textbf{For economists:} MSE loss $\Leftrightarrow$ OLS objective. Cross-entropy $\Leftrightarrow$ multinomial logit likelihood.
\end{itemize}
\end{frame}

\begin{frame}{Training Terminology}
\begin{itemize}
\item \textbf{Epoch:} One complete pass through the entire training dataset.

\medskip
\item \textbf{Mini-batch:} A small subset of training data used for one gradient update.

\medskip
\item \textbf{Iteration:} One parameter update (= one mini-batch processed).
\begin{itemize}
\item If dataset has 10,000 samples and batch size is 100, one epoch = 100 iterations.
\end{itemize}

\medskip
\item \textbf{Automatic Differentiation:} Computational technique for exact gradient computation --- the engine behind \texttt{loss.backward()} in PyTorch.

\medskip
\item \textbf{Episode:} (Reinforcement learning context) One sequence of steps from initial state to terminal state --- analogous to one simulated life-cycle path.
\end{itemize}
\end{frame}

\begin{frame}{Overfitting and Stopping Criteria}
\begin{itemize}
\item \textbf{The fundamental tension:}
\begin{itemize}
\item Training loss: always decreasing with more epochs
\item Validation loss: decreases, then \textbf{increases} $\to$ overfitting
\end{itemize}

\medskip
\item \textbf{Stopping criteria:}
\begin{enumerate}
\item Fixed number of epochs (simple but suboptimal)
\item Monitor validation loss --- stop when it stops improving
\item \textbf{Early stopping:} halt training when validation loss increases for $k$ consecutive epochs
\end{enumerate}

\medskip
\item \textbf{Model selection:} Choose the model (epoch checkpoint) with best validation performance.

\medskip
\item \textbf{Economist's analogy:} Overfitting is the ML version of data mining --- the model memorizes noise in the training sample rather than learning the true data-generating process. Out-of-sample validation is the discipline.
\end{itemize}
\end{frame}

\begin{frame}{Local Optima vs.\ Global Optimum}
\begin{itemize}
\item \textbf{Non-convex loss landscape:}
\begin{itemize}
\item Neural network loss functions are \emph{not} convex --- many local optima and saddle points
\item Global optimum exists but is generally unreachable in practice
\end{itemize}

\medskip
\item \textbf{Mitigation strategies:}
\begin{itemize}
\item Random initialization of weights (different starting points)
\item Momentum and adaptive learning rates (escape shallow minima)
\item Mini-batch stochasticity (natural exploration of the loss surface)
\end{itemize}

\medskip
\item \textbf{Practical reality:} Finding the global optimum is neither necessary nor desirable. A \emph{good local optimum} that generalizes well is sufficient.

\medskip
\item \textbf{Research frontier:} ``Loss landscape'' analysis shows that in over-parameterized networks, most local minima achieve similar performance --- the landscape is ``benign'' (Choromanska et al., 2015).
\end{itemize}
\end{frame}

\begin{frame}{Regularization: Preventing Overfitting}
\begin{itemize}
\item \textbf{L1 / L2 Regularization (weight decay):}
\[
\mathcal{L}_{\text{reg}} = \mathcal{L}_{\text{data}} + \lambda \|\theta\|_p, \quad p \in \{1, 2\}
\]
\begin{itemize}
\item L1 ($p=1$): promotes sparsity (like LASSO in econometrics)
\item L2 ($p=2$): shrinks weights toward zero (like Ridge regression)
\end{itemize}

\medskip
\item \textbf{Dropout:} Randomly set a fraction of neurons to zero during each training step.
\begin{itemize}
\item Forces the network to learn redundant representations
\item Equivalent to training an ensemble of sub-networks
\end{itemize}

\medskip
\item \textbf{Batch Normalization:} Normalize activations within each layer across the mini-batch. Stabilizes and accelerates training.

\medskip
\item \textbf{Data Augmentation:} Expand the training set with transformed copies (rotation, scaling, noise injection).
\end{itemize}
\end{frame}

%=============================================================================
% SECTION 7: SPECIALIZED NEURAL NETWORKS
%=============================================================================
\section{Specialized Architectures}

\begin{frame}{Convolutional Neural Networks (CNNs): Learned Manifold Maps}
\begin{itemize}
\item Recall the two-stage model: $\mathbf{x} \xrightarrow{\;\pi\;} \mathbf{z} \xrightarrow{\;g_\phi\;} \hat{\mathbf{p}}$

\medskip
\item \textbf{In a CNN:}
\begin{itemize}
\item Stacked \textbf{convolution + nonlinearity + pooling} layers implement a \emph{learned} $\pi$
\item Each filter detects a local pattern (edges, curves, textures)
\item \textbf{Convolution:} $Y(i,j) = \sum_m \sum_n X(i-m, j-n)\, W(m,n)$
\item \textbf{Pooling:} $\text{MaxPool}(i,j) = \max_{(m,n) \in R} X(m,n)$ --- spatial summarization
\end{itemize}

\medskip
\item The fully connected ``classifier head'' is the $g_\phi$ map turning features into predictions.

\medskip
\item \textbf{Economic analogy:} CNN layers learn a \emph{sufficient statistic} for prediction --- like finding a low-dimensional state in a high-dimensional macro model.
\end{itemize}
\end{frame}

\begin{frame}{AlexNet: A Landmark CNN Architecture}
\begin{center}
\includegraphics[width=15cm]{alexnet}
\end{center}

{\small Krizhevsky, Sutskever, Hinton (2012). Won ImageNet with 15.3\% top-5 error (vs.\ 26.2\% for non-DL methods). This result launched the deep learning revolution.}
\end{frame}

\begin{frame}{Input Convex Neural Networks (ICNNs)}
\begin{itemize}
\item \textbf{Motivation:} In economics, many functions are known to be convex or concave: cost functions, duality mappings, value functions (after sign change).

\medskip
\item \textbf{ICNN architecture} (Amos, Xu, Kolter, 2017): ensure $f(x;\theta)$ is convex in $x$ by construction:
\begin{align*}
\mathbf{z}^{(0)} &= x \\
\mathbf{z}^{(\ell+1)} &= \sigma\!\left(\underbrace{W^{(\ell)}}_{\geq 0} \mathbf{z}^{(\ell)} + A^{(\ell)} x + b^{(\ell)}\right), \quad \ell = 1,\ldots,L-1
\end{align*}

\item \textbf{Key constraints:}
\begin{itemize}
\item Weights $W^{(\ell)} \geq 0$ (element-wise non-negative) for $\ell \geq 1$
\item Activation $\sigma$ must be convex and non-decreasing (ReLU, softplus)
\item ``Pass-through'' connections $A^{(\ell)} x$ can have unconstrained weights
\end{itemize}

\medskip
\item \textbf{Result:} The composition of convex, non-decreasing functions with non-negative weights preserves convexity.

\medskip
\item \textbf{For concavity:} Use $g(x) = -f(x;\theta)$ where $f$ is an ICNN.
\end{itemize}
\end{frame}

\begin{frame}{Monotone Neural Networks}
\begin{itemize}
\item \textbf{Motivation:} Economic theory often implies monotonicity:
\begin{itemize}
\item Demand is non-increasing in price (law of demand)
\item Consumption is non-decreasing in wealth (normal goods)
\item Value functions are non-decreasing in the state
\end{itemize}

\medskip
\item \textbf{Architecture for monotonicity in input $x_j$:}
\begin{itemize}
\item Constrain all weights connecting from $x_j$ through the network to be \textbf{non-negative}
\item Use \textbf{non-decreasing} activation functions (ReLU, sigmoid, softplus)
\item Composition of non-decreasing functions is non-decreasing
\end{itemize}

\medskip
\item \textbf{Partial monotonicity:} Can enforce monotonicity w.r.t.\ \emph{some} inputs while leaving others unconstrained --- matching economic theory that demand is monotone in price but non-monotone in income (Giffen goods aside).

\medskip
\item \textbf{Practical enforcement:} After each gradient step, project weights to $\mathbb{R}_+$ via $w \leftarrow \max(w, 0)$ or use $w = \exp(\tilde{w})$ reparameterization.
\end{itemize}
\end{frame}

\begin{frame}{Economic Applications of Constrained Architectures}
\begin{itemize}
\item \textbf{ICNN for optimal transport / Wasserstein distance:}
\begin{itemize}
\item Kantorovich duality: $W_1(\mu,\nu) = \sup_{f \text{ 1-Lipschitz}} \mathbb{E}_\mu[f] - \mathbb{E}_\nu[f]$
\item ICNNs parameterize the convex conjugate --- used in generative models and distributional analysis
\end{itemize}

\medskip
\item \textbf{Concave value functions:}
\begin{itemize}
\item In many macro models, $V(k)$ is concave in capital
\item Approximate $V$ with $-\text{ICNN}(k)$ to enforce concavity \emph{by construction}
\item Eliminates spurious non-concavities from unconstrained approximation
\end{itemize}

\medskip
\item \textbf{Monotone policy functions:}
\begin{itemize}
\item Savings policy $a'(a,z)$ is typically increasing in current assets $a$
\item Monotone networks guarantee this without post-hoc corrections
\end{itemize}

\medskip
\item \textbf{General principle:} Encoding economic structure into the network architecture reduces the hypothesis space, improves sample efficiency, and produces economically interpretable solutions.
\end{itemize}
\end{frame}

%=============================================================================
% SECTION 8: CASE STUDY AND PYTORCH
%=============================================================================
\section{Case Study: Curve Fitting and PyTorch}

\begin{frame}{Case Study: DNN vs.\ Polynomial vs.\ Chebyshev}
\begin{itemize}
\item \textbf{Objective:} Approximate a ``kinked'' function using three methods:
\begin{itemize}
    \item Deep Neural Network (DNN)
    \item Traditional polynomial approximation (\texttt{np.polyfit})
    \item Chebyshev polynomial approximation
\end{itemize}

\medskip
\item \textbf{Why this matters for economists:}
\begin{itemize}
    \item Many economic models feature kinks: borrowing constraints, tax brackets, occasionally binding constraints
    \item Polynomials suffer from \textbf{Gibbs phenomenon} near discontinuities
    \item DNNs with ReLU activations handle kinks naturally (piecewise linear)
\end{itemize}

\medskip
\item \textbf{Test function:}
\[
y(x) = \begin{cases} \alpha\beta\, x^\alpha & x \leq 0.5 \\ 2\alpha\beta\, x^\alpha & x > 0.5 \end{cases}
\]
A power function with a discrete jump at $x=0.5$ --- like a policy threshold.
\end{itemize}
\end{frame}

\begin{frame}{Example: GDP Forecasting with a Tiny MLP}
\begin{itemize}
\item \textbf{Objective:} Predict next-quarter real GDP growth $g_t$ from four lags.

\item \textbf{Input:} $\mathbf{x}_t = (g_{t-1}, g_{t-2}, g_{t-3}, g_{t-4}) \in \mathbb{R}^4$

\item \textbf{Architecture:}
\[
\text{Input (4)} \xrightarrow{\text{ReLU}} \text{Hidden (16)} \xrightarrow{\text{ReLU}} \text{Hidden (8)} \xrightarrow{\text{Linear}} \text{Output (1)}
\]

\item \textbf{Why a neural network?}
\begin{itemize}
    \item Standard AR models assume linear persistence
    \item GDP dynamics are often \emph{state-dependent}: shocks have different persistence in recessions vs.\ expansions
    \item The MLP allows coefficients to vary non-linearly with the lag structure
\end{itemize}

\item \textbf{Loss function:} MSE $\equiv$ Gaussian MLE:
$\mathcal{L}(\theta) = \frac{1}{N}\sum_t (\hat{g}_t - g_t)^2$
\end{itemize}
\end{frame}

\begin{frame}[fragile]{PyTorch: Defining and Training the Model}

\begin{lstlisting}[style=pythonstyle]
import torch
import torch.nn as nn
import torch.optim as optim

class GDPNet(nn.Module):
    def __init__(self):
        super().__init__()
        self.net = nn.Sequential(
            nn.Linear(4, 16), nn.ReLU(),
            nn.Linear(16, 8), nn.ReLU(),
            nn.Linear(8, 1)
        )
    def forward(self, x):
        return self.net(x)

model = GDPNet()
criterion = nn.MSELoss()
optimizer = optim.Adam(model.parameters(), lr=0.001)

for epoch in range(1000):
    optimizer.zero_grad()          # 1. Reset gradients
    y_hat = model(X_batch)         # 2. Forward pass
    loss = criterion(y_hat, y_batch) # 3. Compute loss
    loss.backward()                # 4. Backpropagation
    optimizer.step()               # 5. Update weights
\end{lstlisting}
\end{frame}

\begin{frame}[fragile]{PyTorch: Class-Based Network with Batch Normalization}

\begin{lstlisting}[style=pythonstyle]
class KinkNet(nn.Module):
    """Network for approximating kinked functions."""
    def __init__(self, input_dim, hidden_dim1,
                 hidden_dim2, output_dim=1):
        super().__init__()
        self.net = nn.Sequential(
            nn.Linear(input_dim, hidden_dim1),
            nn.BatchNorm1d(hidden_dim1),
            nn.ReLU(),
            nn.Linear(hidden_dim1, hidden_dim2),
            nn.BatchNorm1d(hidden_dim2),
            nn.ReLU(),
            nn.Linear(hidden_dim2, output_dim)
        )

    def forward(self, x):
        return self.net(x)
\end{lstlisting}

\medskip
\textbf{Key additions:} \texttt{BatchNorm1d} normalizes activations within each mini-batch --- stabilizes training and acts as light regularization.
\end{frame}

%=============================================================================
% SECTION 9: LAB PREVIEW
%=============================================================================
\section{Lab Preview}

\begin{frame}{Lab: Machine Learning Basics and Applications}

\begin{tcolorbox}[colback=green!3, colframe=green!40, title=Lab 5A: Machine Learning Basics]
\textbf{Notebook:} \texttt{Lab5A\_MachineLearning\_Basics.ipynb}
\begin{itemize}
\item Build and train neural networks from scratch in PyTorch
\item Compare DNN, polynomial, and Chebyshev approximation on kinked functions
\item Visualize the training process: loss curves, learned function shapes
\item Experiment with architecture choices: depth, width, activation functions
\end{itemize}
\end{tcolorbox}

\medskip

\begin{tcolorbox}[colback=blue!3, colframe=blue!40, title=Lab 5B: FOMC Minutes and 10-Year Yield Model]
\textbf{Notebook:} \texttt{Lab5B\_fomc\_minutes\_10y\_model.ipynb}
\begin{itemize}
\item Apply ML to a real economic forecasting task
\item Extract features from FOMC meeting minutes (text $\to$ numerical features)
\item Train an MLP to predict 10-year Treasury yield movements
\item Evaluate: out-of-sample performance, feature importance
\end{itemize}
\end{tcolorbox}
\end{frame}

%=============================================================================
% SUMMARY
%=============================================================================

\begin{frame}{Summary}
\begin{enumerate}
\item \textbf{Why ML:} Unstructured data, high-dimensional models, and non-linear dynamics make ML essential for modern macro research.

\medskip
\item \textbf{Neural Networks:} Universal approximators that compose linear transformations with nonlinear activations. Barron (1993): break the curse of dimensionality.

\medskip
\item \textbf{Training:} Forward pass $\to$ loss $\to$ backpropagation $\to$ gradient descent. Mini-batch SGD with Adam is the practical standard.

\medskip
\item \textbf{Regularization:} Weight decay, dropout, early stopping, batch normalization prevent overfitting --- the ML analogue of avoiding data mining.

\medskip
\item \textbf{Constrained Architectures:} ICNNs enforce convexity; monotone networks enforce monotonicity. Encoding economic theory into network design improves both accuracy and interpretability.

\medskip
\item \textbf{Next:} Lec.\ 4 applies these tools to \emph{solve} macro models --- deep learning meets the Bellman equation.
\end{enumerate}
\end{frame}

\end{document}
